\documentclass[11pt,a4paper]{article}
\usepackage[margin=23mm]{geometry}
\usepackage{amsmath,amssymb,booktabs,microtype}
\usepackage[hidelinks]{hyperref}
\newcommand{\id}{\mathbb I}
\newcommand{\dd}{\mathrm d}
\title{Route G-R4: A Physical Source and the\protect\\Cross-Transition Universality Test}
\author{Thermal electric-dipole response, without a fitted time law}
\date{24 September 2026}
\begin{document}
\maketitle
\begin{abstract}
We select one physical source whose coupling is known: an isotropic
blackbody radiation bath acting on two clock transitions of the same
$^{171}\mathrm{Yb}^+$ ion. Independently measured polarizabilities fix
both frequency-shift signs. The leading fractional responses differ
by a factor of about $7.25$, excluding a sole common rescaling as a
description of this selected thermal channel. An additional exactly
common multiplicative factor is invisible to the same-bath frequency
ratio and is not excluded. This is a conventional electromagnetic
control, not a derivation of emergent time or an identification of
radiation density with the source document's $E_d$.
\end{abstract}

\section{One declared source, two independently calibrated responses}
Retain the G-R3 lesson: a density observable does not fix a coupling
or its sign. Here the source is thermal electromagnetic radiation,
the clock is one ion, and the microscopic interaction is
\begin{equation}\label{eq:interaction}
 H_{\rm int}=-\mathbf d\cdot\mathbf E(\mathbf x_{\rm ion}).
\end{equation}
The two transitions share the $^2S_{1/2}(F=0)$ ground level $g$:
E2 ends at $^2D_{3/2}(F=2)$ and E3 at $^2F_{7/2}(F=3)$.
The labels E2/E3 describe the \emph{clock interrogation} multipoles;
the selected thermal shift is an electric-\emph{dipole} response of
their levels. An isotropic bath removes the tensor contribution.

Use the single sign convention $\Delta\alpha_i=\alpha_{e_i}-\alpha_g$.
The independent applied-field determinations used here are
\begin{equation}\label{eq:inputs}
 \Delta\alpha_2(0)=6.9(1.4)\,10^{-40},\qquad
 \Delta\alpha_3(0)=0.888(0.016)\,10^{-40}
 \quad [\mathrm{J\,m^2/V^2}].
\end{equation}
The E2 source~\cite{schneider} reports the \emph{opposite} convention,
$\alpha_g-\alpha_{e_2}=-6.9(1.4)\,10^{-40}$; E3~\cite{huntemann}
uses excited minus ground. These are not coefficients adjusted to
our proposed temperature comparison, nor claimed to be the latest
possible calibration. Their uncertainties and provenance are retained.

\textbf{Scope boundary.} This realization assumes ordinary space,
QED and weak-field atomic spectroscopy. Its $u_\gamma$ is thermal
energy per \emph{physical volume}, not G-R3 excitation energy per
addressed cell and not an operational definition of Chen's $E_d$.
It does not complete G-R2's recording apparatus or replace the finite
G-R3 stationary model. It tests a proposed physical interpretation.

\newpage
\section{Derive the thermal response and its sign}
For a monochromatic field with peak amplitude $E_0$ and polarization
$\boldsymbol\epsilon$, second-order perturbation theory gives
\begin{align}
 \delta E_a(\omega)
 &=\frac{E_0^2}{4}\sum_b |\langle b|\mathbf d\cdot
 \boldsymbol\epsilon|a\rangle|^2
 \left[\frac{1}{E_a-E_b+\hbar\omega}
       +\frac{1}{E_a-E_b-\hbar\omega}\right]\\
 &=-\frac14\alpha_a(\omega)E_0^2
  =-\frac12\alpha_a(\omega)\langle E^2\rangle .\label{eq:stark}
\end{align}
Principal values select the dispersive real shift; resonant excitation,
linewidths and scattering are not included in this expression.
After isotropic polarization averaging, a fixed-state scalar response is
\begin{equation}\label{eq:alpha}
 \alpha_a^{\rm sc}(\omega)=\frac23\sum_b
 \frac{(E_b-E_a)\sum_{j=x,y,z}|\langle b|d_j|a\rangle|^2}
 {(E_b-E_a)^2-(\hbar\omega)^2} .
\end{equation}
The sum includes the relevant substates. Upward and downward virtual
transitions contribute with different signs. Even if each individual
level polarizability were positive, its \emph{difference} need not be.
Thus positivity of energy density alone cannot fix a clock-shift sign.
Equation~\eqref{eq:inputs} supplies the independently measured sign.

For the two-polarization Planck distribution, excluding vacuum energy,
\begin{align}
 u_\omega(T)&=\frac{\hbar\omega^3}{\pi^2c^3}
       \frac{1}{e^{\hbar\omega/(k_BT)}-1},\qquad
 u_\gamma=\int_0^\infty u_\omega\dd\omega=a_{\rm rad}T^4,\\
 a_{\rm rad}&=\frac{8\pi^5k_B^4}{15h^3c^3},\qquad
 \langle E^2\rangle=\frac{u_\gamma}{\epsilon_0}.\label{eq:planck}
\end{align}
The integral follows from $\int_0^\infty x^3/(e^x-1)\dd x=\pi^4/15$.
Electric and magnetic fields each carry half the thermal energy;
using $2u_\gamma/\epsilon_0$ for the mean squared electric field would
double-count it. Combining the level shifts yields
\begin{align}
 \delta\nu_i(T)&=-\frac{1}{2h\epsilon_0}\,\mathcal P
       \int_0^\infty\Delta\alpha_i(\omega)u_\omega(T)\dd\omega,
       \label{eq:full}\\
 s_i(T):=\frac{\delta\nu_i}{\nu_i^0}
       &=\beta_i u_\gamma[1+\eta_i(T)],\qquad
 \beta_i=-\frac{\Delta\alpha_i(0)}{2h\epsilon_0\nu_i^0}.\label{eq:beta}
\end{align}
Here $1+\eta_i$ is the spectral integral divided by
$\Delta\alpha_i(0)u_\gamma$, not a fitted slowdown function.
The leading static approximation sets $\eta_i=0$. Positive differential
polarizabilities give negative shifts for both chosen transitions.
This is the conventional BBR shift formula, also used in the 2024
clock analysis~\cite{tofful}.

\newpage
\section{What equal fractional shifts would actually require}
\textbf{Diagonal-level criterion.} Let $H_0$ have level energies $E_a$
and let $\delta E_a$ be their secular first-order shifts. On a connected
graph of measured transitions, every edge has the same fractional
shift $f$ if and only if
\begin{equation}\label{eq:lemma}
 \delta E_a=fE_a+C\quad\hbox{on those levels},\qquad
 \delta H_{\rm sec}=fH_0+C\id .
\end{equation}
\emph{Proof.} On an edge $(a,b)$, the hypothesis is
$\delta E_a-\delta E_b=f(E_a-E_b)$, so
$\delta E_a-fE_a=\delta E_b-fE_b$. Connectivity makes this value one
constant $C$. Conversely that form gives the stated shift on every
edge. This is a statement about diagonal energy shifts, not an
operator identity for arbitrary off-diagonal perturbations or a proof
of universality for untested transitions.

For the three levels $g,e_2,e_3$, the selected leading thermal response
would therefore require
\begin{equation}\label{eq:necessary}
 \frac{\Delta\alpha_2(0)}{\nu_2^0}
 =\frac{\Delta\alpha_3(0)}{\nu_3^0},\qquad \beta_2=\beta_3 .
\end{equation}
The independent inputs violate this condition. In the off-resonant
low-temperature expansion, analytic $\Delta\alpha_i(\omega)$ gives
$s_i=\beta_i a_{\rm rad}T^4+O(T^6)$. Different leading coefficients
cannot be made identical on a low-temperature interval by $T^6$
terms. This asymptotic observation is not a full finite-temperature
error bound or an exclusion of accidental equality at one temperature.

Define $R(u)=\nu_3(u)/\nu_2(u)$ and $R_0=\nu_3^0/\nu_2^0$.
Within the retained static-response model, exact quotient algebra gives
\begin{align}
 \frac{R(u)}{R_0}-1
 &=\frac{(\beta_3-\beta_2)u}{1+\beta_2u},\label{eq:ratio}\\
 \frac{R(u_h)}{R(u_l)}-1
 &=\frac{(\beta_3-\beta_2)(u_h-u_l)}
 {(1+\beta_2u_h)(1+\beta_3u_l)}.\label{eq:modulation}
\end{align}
Keeping these denominators is convenient algebra, not control of all
higher orders in the microscopic interaction. The expressions avoid
subtracting nearly equal floating-point numbers.

\textbf{Identifiability limit.} Suppose an additional common physical
factor exists:
\begin{equation}\label{eq:blind}
 \nu_i(u)=F(u)\nu_i^0[1+s_i(u)],\qquad
 R(u)=R_0\frac{1+s_3(u)}{1+s_2(u)} .
\end{equation}
$F$ cancels exactly, even when it changes between setpoints. A
nonconstant ratio rejects a \emph{sole} common factor as the explanation
of this thermal response; it does not exclude a common factor in
addition to it. Equal fractions would be necessary, not sufficient,
evidence for a universal time effect. Relative phases obey
$\dd\Phi_3/\dd\Phi_2=\nu_3/\nu_2$ in a quasistatic environment,
without needing a preferred external time unit for that comparison.

\newpage
\section{Pinned numerical prediction and its uncertainty boundary}
Use rounded frequencies $\nu_2^0=688.359\,\mathrm{THz}$ and
$\nu_3^0=642.121\,\mathrm{THz}$, exact SI $h,k_B,c$, and rounded
$\epsilon_0=8.8541878\,10^{-12}\,\mathrm{F/m}$.
These suffice for a response-structure test, not a metrological
absolute-frequency recommendation. At $T=300\,\mathrm K$,
\begin{equation}
 u_\gamma=6.12824\,10^{-6}\,\mathrm{J/m^3},\qquad
 E_{\rm rms}=831.943\,\mathrm{V/m}.
\end{equation}
The static electric-dipole predictions ($\eta_2=\eta_3=0$) are:
\begin{center}
\begin{tabular}{@{}lrr@{}}
\toprule
Transition & $\delta\nu\ [\mathrm{Hz}]$ & $s=\delta\nu/\nu^0$\\
\midrule
E2 & $-0.3604$ & $-5.235\,10^{-16}$\\
E3 & $-0.04638$ & $-7.223\,10^{-17}$\\
\bottomrule
\end{tabular}
\end{center}
Thus $|\beta_2/\beta_3|\simeq7.25$, and
\begin{equation}
 \frac{R(300\,\mathrm K)}{R(0)}-1\simeq4.513\,10^{-16},\qquad
 \frac{R(310\,\mathrm K)}{R(300\,\mathrm K)}-1\simeq6.325\,10^{-17}.
\end{equation}
Both shifts are downward, but E2 falls by a larger fraction, so E3/E2
increases. The second number is a two-setpoint \emph{prediction}, not
an observed temperature-cycle signal. Digits document reproducibility;
the polarizability uncertainty, especially E2's, sets physical precision.

The independently published E3 correction
$\eta_3(300\,\mathrm K)=-0.0015(7)$~\cite{huntemann} changes its
shift to about $-0.04631\,\mathrm{Hz}$. We do not supply a complete
E2 spectral response, full covariance or apparatus radiation model,
and do not label the table as a full finite-temperature prediction.

For a transparent robustness check, allow each polarizability to vary
over its separate $\pm3\sigma$ input interval and each dynamic factor
over $[0.9,1.1]$. Monotonicity gives the exact box lower bound
\begin{equation}
 (\beta_3-\beta_2)_{\min}
 =\frac{1}{2h\epsilon_0}
 \left[\frac{0.9(\alpha_2-3\sigma_2)}{\nu_2^0}
       -\frac{1.1(\alpha_3+3\sigma_3)}{\nu_3^0}\right]
 =1.642\,10^{-11}\,\mathrm{m^3/J}>0,
\end{equation}
where $\alpha_i$ denotes $\Delta\alpha_i(0)$ in this bound.
The corresponding leading ratio offset at 300 K exceeds
$1.006\,10^{-16}$. This is a \emph{conditional deterministic stress box},
not a statistical confidence level or a derived bound on all missing
dynamic corrections. The unknown cross-paper correlation $\rho$ gives
\begin{equation}
 \sigma_{\beta_3-\beta_2}^2=\sigma_{\beta_3}^2+\sigma_{\beta_2}^2
              -2\rho\sigma_{\beta_3}\sigma_{\beta_2},\quad
 \sigma_{\beta_3-\beta_2}\in
 [|\sigma_{\beta_2}-\sigma_{\beta_3}|,\sigma_{\beta_2}+\sigma_{\beta_3}].
\end{equation}
No independence-based significance or new combined best value is claimed.

\newpage
\section{The next physical discriminator, not another time-law fit}
\textbf{A: test the differential prediction.} Specify a real common
radiation environment and infer its spectrum at the ion, not merely
a remote wall temperature. Interleave E2/E3 interrogation at two
stable setpoints and retain the raw ratio before subtracting the
deliberately modulated BBR response. Same-ion interleaved comparisons
are an established protocol~\cite{filzinger}, but no new data are
analyzed here. Published BBR-corrected clock outputs are not independent
temperature-response observations. Returning a flat ratio after
subtracting an assumed BBR law would not validate that same law.

The temperature change may also alter tensor shifts, magnetic fields,
trap motion, probe shifts and thermal gradients. These must be measured
or bounded. Magnetic-dipole/quadrupole thermal contributions and
linewidth effects are outside our selected scalar electric-dipole
calculation; they are not asserted to vanish in a real apparatus.
The sign and leading unequal responses are a completed source audit;
experimental closure requires those controls, not more Gaussian scans.

\textbf{B: make a common factor observable.} If the goal is specifically
$F(u)$ in Eq.~\eqref{eq:blind}, provide an independently specified
reference outside the changed bath, with a physical signal-comparison
protocol and its own response model, or an independently predicted
non-clock observable. Do not assume by definition that the reference
is unaffected. A factor that changes all physical standards identically
may remain a units/reparametrization freedom rather than an observable.

\textbf{C: density versus spectral content.} As a separate optional test,
compare two nonthermal radiation spectra with the same integrated
$u_\gamma$ but different spectral weights. Equation~\eqref{eq:full}
predicts different shifts when $\Delta\alpha_i(\omega)$ varies.
This directly tests whether scalar density is sufficient; it requires
spectral calibration and is not a new free per-clock parameter.

\textbf{Completed deliverable.} G-R4 supplies a source, an independently
fixed sign, the cross-transition criterion and a reproducible numerical
prediction. The preferred theoretical revision is an interaction-derived
\emph{spectral response}, with a universal component treated separately
and only when observable. No additional $F(u)$ is fitted. The code's
91 checks validate the displayed algebra and implementation, not a
new empirical time effect. Input conventions, hashes, limitations and
roadmap are stored alongside the TeX; earlier Route-G results are preserved.

\begin{thebibliography}{9}\small
\bibitem{schneider} T. Schneider, E. Peik and Chr. Tamm,
\emph{Sub-Hertz Optical Frequency Comparisons between Two Trapped
$^{171}$Yb$^+$ Ions}, Phys. Rev. Lett. \textbf{94}, 230801 (2005),
\href{https://doi.org/10.1103/PhysRevLett.94.230801}{doi:10.1103/PhysRevLett.94.230801}.
\bibitem{huntemann} N. Huntemann et al., \emph{Single-Ion Atomic Clock
with $3\times10^{-18}$ Systematic Uncertainty}, Phys. Rev. Lett.
\textbf{116}, 063001 (2016),
\href{https://doi.org/10.1103/PhysRevLett.116.063001}{doi:10.1103/PhysRevLett.116.063001}.
\bibitem{tofful} A. Tofful et al., \emph{$^{171}$Yb$^+$ optical clock
with $2.2\times10^{-18}$ systematic uncertainty and absolute frequency
measurements}, Metrologia \textbf{61}, 045001 (2024),
\href{https://doi.org/10.1088/1681-7575/ad53cd}{doi:10.1088/1681-7575/ad53cd}.
\bibitem{filzinger} M. Filzinger et al., \emph{Improved limits on the
coupling of ultralight bosonic dark matter to photons from optical
atomic clock comparisons} (2023),
\href{https://arxiv.org/abs/2301.03433}{arXiv:2301.03433}.
\end{thebibliography}
\end{document}
